\section{Wstęp}
\label{sec:introduction}

\TODO{problem: modele językowe działają jako agenci przez wywołania narzędzi, a testy tool callingu oceniają głównie
poprawność wywołań i wyniku, rzadko to, jak model doszedł do rozwiązania.}

\subsection{Cel badania}
\label{sec:goal}

\TODO{cel: otwarte środowisko testowe z analizatorem przebiegów i raport, jak skutecznie i jakimi algorytmami modele
LLM różnych dostawców rozwiązują zadania przez tool calling.}

\subsection{Pytania badawcze}
\label{sec:questions}

\begin{enumerate}
  \item \textbf{Skuteczność:} jak odsetek poprawnie rozwiązanych zadań zależy od modelu, dostawcy i rozmiaru
    zadania $N$?
  \item \textbf{Algorytm:} jakie algorytmy modele faktycznie wykonują: znane, ich mieszanki czy własne strategie?
  \item \textbf{Deklaracja a wykonanie:} czy algorytm zapowiedziany w tekście zgadza się z wykonanym?
  \item \textbf{Efektywność:} ile porównań i zamian zużywa model wobec granic teoretycznych i algorytmów wzorcowych?
  \item \textbf{Zmiana strategii:} czy model zmienia algorytm, gdy rośnie $N$ albo gdy narzędzie zwróci błąd?
\end{enumerate}

\subsection{Wkład}
\label{sec:contribution}

\TODO{co nowego: każdy ruch jako osobne wywołanie narzędzia na ukrytym stanie, rozpoznawanie algorytmu z zapisu
ruchów, porównanie deklaracji z wykonaniem, gotowe modele wielu dostawców bez trenowania.}
